\section{El reto}\label{sec:reto}
Las 59 parcelas a predecir forman el conjunto de \emph{predicción} y las 138 con rendimiento conocido el de
\emph{entrenamiento}. Todas son del ciclo de abril a octubre de 2025~\cite{bases}. El objetivo es el rendimiento en
toneladas por hectárea (t/ha), no la calidad del grano. Son variables distintas y a veces opuestas, porque más
nitrógeno suele subir el rendimiento pero también la proteína, y una proteína alta empeora la cebada para malta.

\subsection{Entregables}
La primera etapa pide seis entregables.
\begin{enumerate}
  \item Las predicciones de rendimiento de las 59 parcelas.
  \item Un reporte técnico con secciones y extensión fijas. Los anexos pueden ocupar hasta cuatro cuartillas y deben incluir una tabla con las predicciones de las 59 parcelas~\cite{preguntas}.
  \item Una aplicación web que ejecute el modelo y permita consultar los resultados.
  \item El código fuente con instrucciones para ejecutarlo.
  \item Un video de máximo cinco minutos.
  \item El registro del uso de inteligencia artificial, con los \emph{prompts} anexos.
\end{enumerate}
Además de la carga en la plataforma, cada equipo envía al correo innovacion@fira.gob.mx un archivo CSV con las predicciones, llamado nombre\_del\_equipo\_prediccion.csv. En el mismo correo se pueden agregar, de forma opcional y breve, enlaces al código o a otros anexos~\cite{preguntas}.

\subsection{Evaluación}
El desempeño predictivo y su justificación técnica valen el 70\% de la primera etapa, y la aplicación el 30\%
(Tabla~\ref{tab:evaluacion}). El primer filtro es el error de predicción, que los organizadores medirán con el error cuadrático medio, y después se valoran la metodología y la aplicación~\cite{preguntas}. Cada equipo debe reportar también su propio error. Hay puntos extra por responder con evidencia tres preguntas guía, qué índice satelital
es el más importante, qué variable climática pesa más y qué combinación de variables funciona mejor. Los tres mejores
equipos presentan su trabajo ante el jurado en una segunda etapa.

\begin{table}[!t]
\caption{Criterios de evaluación}\label{tab:evaluacion}
\centering\footnotesize
\begin{tabularx}{\columnwidth}{@{}lXr@{}}
\toprule
Etapa & Criterio & Peso \\
\midrule
1 & Desempeño predictivo (RMSE, MAE, $R^2$ u otras métricas) y justificación técnica & 70\% \\
1 & Aplicación o interfaz web & 30\% \\
1 & Respuestas a las tres preguntas guía & Extra \\
2 & Presentación de 15 minutos ante el jurado & Define lugares \\
\bottomrule
\end{tabularx}
\end{table}

\subsection{Fechas y reglas}
La entrega cierra el 19 de octubre de 2026 antes de la medianoche. Los finalistas se anuncian el 3 de noviembre, la
presentación es el 10 de noviembre y la premiación el 20 de noviembre. Las bases permiten usar datos públicos si se
documenta la fuente y la forma en que se incorporaron, y permiten usar inteligencia artificial si se indica dónde se usó. Se permite un asesor que no esté inscrito, siempre que se reporte la aportación de cada integrante. La aplicación debe servir tanto a analistas de FIRA como a productores, y el jurado está formado por profesores universitarios~\cite{preguntas}.
